\section{Technical Implementation of the RAG Module}
\label{sec:rag}

This section explains the Retrieval-Augmented Generation (RAG) module in a clear,
defence-ready way: what problem it solves, how it works in four short steps, and how
it connects to the dashboard~\cite{lewis2020rag}.

\subsection{What RAG Is (and Why We Use It)}
\label{sec:rag-why}

A normal chatbot answers from its training memory. That is risky for regulation:
it can invent article numbers, miss new laws, and give answers without sources.

RAG works differently:
\begin{enumerate}[noitemsep]
  \item the user asks a question;
  \item the system \textbf{retrieves} relevant passages from our curated corpus;
  \item the system \textbf{answers only from those passages};
  \item every answer shows \textbf{citations}.
\end{enumerate}

In one sentence: \emph{retrieve first, then generate}. This reduces hallucination and
keeps the student accountable for every claim~\cite{lewis2020rag,karpukhin2020}.

\subsection{Brief RAG Architecture}
\label{sec:rag-arch}

Figure~\ref{fig:rag-arch} shows the architecture in four blocks only.

\begin{figure}[htbp]
\centering
\begin{tikzpicture}[
  box/.style={draw=aivNavy, line width=0.8pt, rounded corners=2pt,
              align=center, font=\small\bfseries, fill=aivBox,
              minimum width=10.6cm, minimum height=0.85cm, inner sep=4pt},
  arr/.style={-{Latex}, thick, aivNavy}]
  \node[box] (a) {1.~Corpus (laws, guidance, literature, country JSON)};
  \node[box, below=0.35cm of a] (b) {2.~Retrieve top-$k$ passages (BM25 + optional dense search)};
  \node[box, below=0.35cm of b] (c) {3.~Grounded answer (only from retrieved text + citations)};
  \node[box, below=0.35cm of c] (d) {4.~Dashboard output (RAG Assistant page / evidence cards)};
  \draw[arr] (a) -- (b);
  \draw[arr] (b) -- (c);
  \draw[arr] (c) -- (d);
\end{tikzpicture}
\caption{Brief RAG architecture used in this project (four steps). Source: author.}
\label{fig:rag-arch}
\end{figure}

\begin{table}[htbp]
\centering
\caption{RAG pipeline in one view}
\label{tab:rag-brief}
\begin{tabular}{@{}clp{8.2cm}@{}}
\toprule
\textbf{Step} & \textbf{Name} & \textbf{What happens} \\
\midrule
1 & Corpus & Store public regulatory texts and curated country narratives as short chunks \\
2 & Retrieve & Rank the most relevant chunks for the user question \\
3 & Generate & Write an answer using only those chunks; add citations \\
4 & Display & Show answer + sources in the React RAG Assistant \\
\bottomrule
\end{tabular}
\end{table}

Formally, for question $q$ and corpus chunks $\mathcal{D}$:
\begin{equation}
\label{eq:rag}
\hat{y}=g\bigl(q,\operatorname{TopK}(q,\mathcal{D})\bigr),
\end{equation}
where $\operatorname{TopK}$ returns the best passages and $g$ produces a grounded answer
with citations. No answer is accepted without retrieved evidence.

\subsection{Corpus and Chunking (Step 1)}
\label{sec:rag-corpus}

\textbf{Inputs.}
Agency documents (FDA, EU AI Act, MDR/IVDR, WHO), peer-reviewed papers, country
narratives in \texttt{compliance\_dataset.json}, and the literature-review markdown.

\textbf{Chunking.}
Long documents are split into short passages (about 350 tokens, overlap 80) so retrieval
can find the exact paragraph. Each chunk keeps simple metadata: source, jurisdiction,
theme tag, and text.

\subsection{Retrieval (Step 2)}
\label{sec:rag-retrieve}

Retrieval finds the best passages for the question:
\begin{itemize}[noitemsep]
  \item \textbf{BM25 (sparse):} good for exact legal terms (e.g.\ ``Article 10'', ``510(k)'');
  \item \textbf{Dense search (optional):} good for similar meaning even if words differ;
  \item \textbf{Hybrid:} combine both (weight $\alpha\approx 0.6$ dense / $0.4$ BM25)~\cite{robertson2009,karpukhin2020}.
\end{itemize}

The dashboard MVP uses an offline lexical retriever over the same corpus interface, so
the defence demo runs without a GPU or paid API. Top-$k$ is typically 6--8 passages.
Optional filters: jurisdiction and compliance theme.

\subsection{Grounded Generation and Citations (Step 3)}
\label{sec:rag-prompt}

The generator must:
\begin{itemize}[noitemsep]
  \item use \emph{only} retrieved passages;
  \item cite chunk identifiers;
  \item say ``insufficient evidence'' when the corpus is not enough;
  \item never invent statute numbers or device counts;
  \item remind the user that output is \textbf{not legal advice}.
\end{itemize}

\begin{lstlisting}[caption={Short grounded prompt used by the RAG module},label={lst:prompt}]
Use ONLY the CONTEXT passages. Cite [chunk_id].
If CONTEXT is insufficient, say "Insufficient evidence".
Do not provide legal advice.
CONTEXT: [1] ... [2] ...
QUESTION: ...
\end{lstlisting}

\subsection{Dashboard Output and HITL (Step 4)}
\label{sec:rag-api}

The React \textbf{RAG Assistant} page shows:
\begin{enumerate}[noitemsep]
  \item the grounded answer;
  \item ranked passages with scores;
  \item links to Country Detail for follow-up.
\end{enumerate}

RAG can suggest evidence for theme scores, but \textbf{humans confirm} any score change
(HITL). Flow:
\[
\text{Sources}\rightarrow\text{Retrieve}\rightarrow\text{Answer + citations}
\rightarrow\text{HITL}\rightarrow\text{JSON / Dashboard}.
\]

\subsection{Simple Defence Example}
\label{sec:rag-demo}

\textbf{Question:} Does the EU AI Act add documentation duties beyond GDPR for high-risk
healthcare AI?

\textbf{Retrieved:} EU AI Act high-risk documentation/logging passages; GDPR health-data
context; EU country transparency narrative.

\textbf{Answer idea:} Yes---the AI Act adds technical documentation and oversight duties
on top of GDPR privacy rules; citations point to the retrieved passages.

\textbf{Next click in UI:} Country Detail $\rightarrow$ European Union $\rightarrow$
transparency / ethics themes.

\subsection{How We Evaluate RAG}
\label{sec:rag-eval}

\begin{table}[htbp]
\centering
\caption{Practical RAG checks for defence}
\label{tab:rag-eval}
\begin{tabular}{@{}p{3.4cm}p{8.0cm}@{}}
\toprule
\textbf{Check} & \textbf{Meaning} \\
\midrule
Relevance & Are the top passages actually about the question? \\
Groundedness & Is each answer sentence supported by a cited chunk? \\
Citation accuracy & Do cited IDs really contain the claim? \\
Abstention & Does the system refuse when evidence is missing? \\
Safety & No invented article numbers in audited answers \\
\bottomrule
\end{tabular}
\end{table}

Baselines compared qualitatively: BM25-only vs hybrid retrieval. Hybrid is preferred
when questions mix exact legal identifiers and semantic wording.

\subsection{Limits (Honest Boundary)}
\label{sec:rag-limits}

\begin{itemize}[noitemsep]
  \item Corpus is mostly English public sources;
  \item Chunking can split long cross-references (mitigated by overlap);
  \item Generation can still over-generalize (mitigated by citations + HITL);
  \item No patient data is indexed---only public regulatory material.
\end{itemize}

\subsection{Summary}
\label{sec:rag-summary}

RAG in this project is a short, clear loop: \textbf{corpus $\rightarrow$ retrieve
$\rightarrow$ grounded answer with citations $\rightarrow$ dashboard}. It is an
information-retrieval tool for compliance research, not a legal chatbot and not a
replacement for expert judgment.
